\begin{resumo}[Abstract]
 \begin{otherlanguage*}{english}
  In agile development, acceptance criteria (AC) state when a user story is done, and when they are written in Gherkin they can be checked automatically. Writing them by hand takes time and domain knowledge, and large language models (LLMs) have started to support this task. It is still unclear, however, whether giving the model the context of the project itself improves the generated criteria. This work aims to evaluate the accuracy of Gherkin scenarios generated by LLMs from user stories, comparing generation with and without retrieval-augmented generation (RAG). In the first stage, a structured literature review covering 25 studies was carried out. The review showed that the topic is recent, with 80\% of the studies published from 2025 onwards, that generation relies mainly on prompt writing, and that the evaluation of generated criteria is still weak: few studies compare them against a reference set, and only two actually follow an empirical strategy. Based on these gaps, an experimental study was planned with a single factor, the use of RAG, applied in a paired design to user stories from the MeasureSoftGram and CAPJu projects. Accuracy will be measured with the Acc@Hit and Acc@Cor metrics, based on the atomic test objectives of the reference set and on the judgment of a committee of three LLMs, and the difference between treatments will be tested with a paired test. Running the study and analyzing its results are planned for the second stage of this work.

  \vspace{\onelineskip}

  \noindent
  \textbf{Keywords}: acceptance criteria; Gherkin; large language models; retrieval-augmented generation; Behavior-Driven Development.
 \end{otherlanguage*}
\end{resumo}
